% ============================================================================
% ARBEITSBLATT DS 3 - Gruppe B: Arbeiten im Ausland (S. 105-106)
% ============================================================================

\documentclass[11pt,a4paper]{article}

\usepackage[utf8]{inputenc}
\usepackage[T1]{fontenc}
\usepackage[ngerman]{babel}
\usepackage[left=2cm,right=2cm,top=1.8cm,bottom=1.5cm]{geometry}
\usepackage{helvet}
\usepackage[table]{xcolor}
\usepackage{tabularx}
\usepackage{array}
\usepackage{enumitem}
\usepackage{tcolorbox}
\usepackage{fancyhdr}
\usepackage{lastpage}

\renewcommand{\familydefault}{\sfdefault}

\definecolor{spanisch}{HTML}{c41e3a}
\definecolor{grau}{HTML}{666666}
\definecolor{hellgrau}{HTML}{f5f5f5}
\definecolor{gruppeB}{HTML}{1565c0}

\pagestyle{fancy}
\fancyhf{}
\fancyhead[L]{\small\textcolor{grau}{Spanisch Spätbeginner -- Gruppe B}}
\fancyhead[R]{\small\textcolor{grau}{AB-03-B}}
\fancyfoot[C]{\small\thepage/\pageref{LastPage}}
\renewcommand{\headrulewidth}{0.5pt}

\newcommand{\luecke}[1]{\underline{\hspace{#1}}}
\newcommand{\punktebox}[1]{\hfill\fbox{\small \_\_/#1}}
\newcommand{\es}[1]{\textit{#1}}

\newtcolorbox{merkkasten}[1][]{
    colback=white,
    colframe=gruppeB,
    fonttitle=\bfseries\color{gruppeB},
    title=#1,
    boxrule=1.5pt,
    arc=0pt,
    left=5pt, right=5pt, top=5pt, bottom=5pt
}

\newtcolorbox{buchverweis}{
    colback=hellgrau,
    colframe=grau,
    boxrule=0.5pt,
    arc=0pt,
    left=5pt, right=5pt, top=3pt, bottom=3pt
}

\newcounter{aufgabe}
\newcommand{\aufgabe}[2]{%
    \stepcounter{aufgabe}%
    \vspace{0.8em}%
    \noindent\textbf{\theaufgabe.} #1 \punktebox{#2}%
    \vspace{0.3em}%
}

\begin{document}

% ============================================================================
% KOPFBEREICH
% ============================================================================
\noindent
\begin{tabularx}{\textwidth}{@{}Xr@{}}
    {\Large\textbf{\textcolor{gruppeB}{Trabajar en el extranjero}}} & Name: \luecke{5cm} \\[0.3em]
    {\small DS 3 -- Arbeiten im Ausland (S. 105--106)} & Datum: \luecke{3cm} \\
\end{tabularx}

\vspace{0.3em}
\hrule
\vspace{0.8em}

% ============================================================================
% MERKKASTEN: Gründe für Auslandsaufenthalt
% ============================================================================
\begin{merkkasten}[Kasten 1: Gründe für einen Auslandsaufenthalt]
\vspace{0.2em}

\begin{tabularx}{\textwidth}{@{}|X|X|@{}}
\hline
\rowcolor{hellgrau}
\textbf{Spanisch} & \textbf{Deutsch} \\
\hline
ganar experiencia & Erfahrung sammeln \\
\hline
aprender idiomas & \luecke{4cm} \\
\hline
conocer otras culturas & \luecke{4cm} \\
\hline
mejorar las oportunidades de trabajo & \luecke{4cm} \\
\hline
desarrollar la independencia & \luecke{4cm} \\
\hline
\end{tabularx}

\vspace{0.3em}
\textbf{Strukturen:} \es{para + Infinitiv} (um zu), \es{porque} (weil), \es{ya que} (da)
\vspace{0.2em}
\end{merkkasten}

\vspace{1em}

% ============================================================================
% AUFGABEN
% ============================================================================

\aufgabe{Ergänze den Kasten oben mit den deutschen Übersetzungen.}{4}

\vspace{0.3em}

\aufgabe{Leseverstehen: Carla und Jan. (Buch S. 105, Aufg. 1)}{8}

\begin{buchverweis}
\textbf{Lies die Texte über Carla und Jan.} Beantworte die Fragen auf Spanisch.
\end{buchverweis}

\vspace{0.3em}
\textbf{Über Carla:}
\begin{enumerate}[label=\alph*), leftmargin=2em]
    \item ¿Dónde trabaja Carla?

    \luecke{12cm}

    \item ¿Por qué decidió ir a Barcelona?

    \luecke{12cm}
\end{enumerate}

\vspace{0.3em}
\textbf{Über Jan:}
\begin{enumerate}[label=\alph*), leftmargin=2em, start=3]
    \item ¿Qué hace Jan en Madrid?

    \luecke{12cm}

    \item ¿Cuáles son las ventajas de su experiencia?

    \luecke{12cm}
\end{enumerate}

\vspace{0.8em}

\aufgabe{Dilo con otras palabras. (Buch S. 106, Aufg. 3a)}{6}

\begin{buchverweis}
\textbf{Erkläre die Ausdrücke mit anderen Worten auf Spanisch.}
\end{buchverweis}

\begin{enumerate}[label=\alph*), leftmargin=2em]
    \item \es{ganar experiencia profesional} =

    \luecke{12cm}

    \item \es{ampliar los conocimientos} =

    \luecke{12cm}

    \item \es{mejorar el currículum} =

    \luecke{12cm}
\end{enumerate}

\vspace{0.8em}

\aufgabe{Wörterbucharbeit: Finde die Bedeutung. (Buch S. 105, Aufg. 2)}{4}

\begin{tabularx}{\textwidth}{@{}|X|X|X|X|@{}}
\hline
\rowcolor{hellgrau}
\textbf{Wort} & \textbf{Bedeutung} & \textbf{Wort} & \textbf{Bedeutung} \\
\hline
las cualificaciones & \luecke{2.5cm} & el empleo & \luecke{2.5cm} \\[0.5em]
\hline
las prácticas & \luecke{2.5cm} & la estancia & \luecke{2.5cm} \\[0.5em]
\hline
\end{tabularx}

\vspace{0.8em}

\aufgabe{Schreibe: Warum würdest du im Ausland arbeiten/studieren wollen?}{10}

\begin{buchverweis}
\textbf{Schreibe 6--8 Sätze.} Verwende: \es{para + Infinitiv}, \es{porque}, \es{ya que}, \es{además}.
\end{buchverweis}

\vspace{0.3em}
\begin{tabularx}{\textwidth}{@{}|X|@{}}
\hline
\\[0.4em]
\luecke{14cm} \\[0.7em]
\luecke{14cm} \\[0.7em]
\luecke{14cm} \\[0.7em]
\luecke{14cm} \\[0.7em]
\luecke{14cm} \\[0.7em]
\luecke{14cm} \\[0.4em]
\hline
\end{tabularx}

\vspace{1em}
\hrule
\vspace{0.3em}
\begin{flushright}
\textbf{Gesamt: \luecke{1cm} / 32 Punkte}
\end{flushright}

\end{document}
